\section{Background}

\subsection{Fuzzing}

Fuzz testing, or \textit{fuzzing} for short, is a widely deployed method for automatic
software testing and vulnerability discovery that makes use of a diverse set of
tools.\cite{8863940}
Its goal is to find inputs to a Program Under Test (PUT) that trigger a certain
behavior, often a crash that might then lead to the discovery of bugs or security
vulnerabilities.
One way to classify fuzzers is by the information that is available to them.

Fuzzers that only make use of externally observable behavior of the PUT are
called \textit{black-box} fuzzers.
In this case, there might only be access to a crash oracle and no further
conclusions about the behavior of the PUT under a certain input can be drawn.
In contrast to this, \textit{white-box} fuzzers allow full access to the PUT,
usually including any original source code or documentation which can greatly
aid the fuzzing process.

Since detailed information is often not available but observing certain parts
of the internal behavior of the PUT might still be possible, \textit{grey-box}
fuzzers represent an important middle-ground between the previous two extremes.
A common source of instrumentation data that grey-box fuzzers make use of is the
\textit{coverage map} of the PUT during execution.
It allows the fuzzer to direct its exploration efforts towards promising regions of the PUT.


One famous example of a coverage-guided greybox fuzzer is AFL++\cite{fioraldi2020afl++}.
It works by iteratively improving the \textit{corpus}, a set of promising
inputs, through a series of mutations and selections.
Mutations include operations like bit flips, substitutions and splicing.
Newly generated inputs are then evaluated and added to the corpus for further
mutation and selection if they produce new coverage.
The success of a greybox fuzzer is strongly dependent on a high-quality initial
corpus and appropriate seed mutation and selection strategies.

\subsection{Reinforcement learning}

Reinforcement learning (RL) is an extremely general framework for approaching problems
in which an agent strives to maximize a reward signal through a selection of actions
in certain states.\cite{sutton1998reinforcement}
These kinds of problems can be modeled as finite Markov decision processes (MDPs).
A Markov decision process is a 3-tuple $MDP=(S,A,p)$ where

\begin{itemize}
  \item $\mathcal{S}$ is the set of states,
  \item $\mathcal{A}$ is the set of actions where $\mathcal{A}_s\subseteq A$ are actions available in state $s\in\mathcal{S}$,
  \item $p(s^\prime,r|s,a)=\Pr\{S_t=s^\prime, R_t=r|S_{t-1}=s, A_{t-1}=a\}$
  describes the system dynamics, i.e., the probability of encountering a next
  state $s^\prime$ and reward $r$ at time step $t$ when taking action $a$ in
  state $s$ at time step $t-1$
\end{itemize}

An RL agent's goal is to learn a policy $\pi:\mathcal{S}\times\mathcal{A}\rightarrow [0,1]$
mapping each state $s$ to probabilities of taking each possible action $a$ to maximize
the cumulative reward.
Algorithms for finding an optimal $\pi^*$ can be divided into three classes:
value-based, policy-based and actor-critic.

Value-based methods try to learn a value function $Q:\mathcal{S}\rightarrow \mathcal{R}$
mapping states to their expected cumulative reward so that in each step, the
action that maximizes the expected reward can be taken.
These include well-known methods such as Q-learning, TD($\lambda$) or SARSA.

Policy-based methods try to directly construct an optimal policy $\pi^*$ without the
use of a value function.
They have various advantages.
For example, they can easily handle continuous action spaces and ensure stochasticity
which is important for state space exploration.
A well-known classical policy-based algorithm is REINFORCE.\cite{williams1992simple}

Actor-critic methods synthesize value- and policy-based methods by including information
about a state's value in the policy learning process.\\

Since classical, tabular methods don't work well with large state spaces,
state-of-the-art methods make use modern-day deep learning advancements.
Examples include the family of Deep-Q Network (DQN)
algorithms\cite{mnih2013playing}\cite{hessel2018rainbow} approximating the $Q$
function by neural networks, policy-based gradient-descent methods like
Proximal Policy Optimization (PPO)\cite{schulman2017proximal} and actor-critic
methods like A3C\cite{mnih2016asynchronous}.
Impressive results range from the master-level Go program
AlphaGo\cite{silver2016mastering} and its more generalized successors
AlphaZero\cite{silver2017mastering} and MuZero\cite{schrittwieser2020mastering}
to Dreamerv3\cite{hafner2023mastering} that is able to learn how to acquire
diamonds in the popular video game \textit{Minecraft} without any human input
or expert knowledge.
